\documentclass{article}
\usepackage[a4paper,margin=25mm]{geometry}
\usepackage{amsmath,amssymb}
\usepackage{xurl}
\usepackage[colorlinks=true,linkcolor=blue,citecolor=blue,urlcolor=blue]{hyperref}
% Shared commands for notes in docs/paper-gaps/.

\DeclareUnicodeCharacter{2080}{\ifmmode{}_{0}\else\textsubscript{0}\fi}
\DeclareUnicodeCharacter{2081}{\ifmmode{}_{1}\else\textsubscript{1}\fi}
\DeclareUnicodeCharacter{2082}{\ifmmode{}_{2}\else\textsubscript{2}\fi}
\DeclareUnicodeCharacter{2099}{\ifmmode{}_{n}\else\textsubscript{n}\fi}

\newcommand{\Lean}{\textsc{Lean}}
\ExplSyntaxOn
\NewDocumentCommand{\leanid}{m}{
  \ifmmode
    \textsf{\detokenize{#1}}
  \else
    \group_begin:
    \urlstyle{sf}
    \tl_set:Nn \l_tmpa_tl {#1}
    \tl_replace_all:Nnn \l_tmpa_tl {\_} {_}
    \exp_args:NV \path \l_tmpa_tl
    \group_end:
  \fi
}
\ExplSyntaxOff
\providecommand{\tr}{\operatorname{tr}}
\newcommand{\tnleanIssueBase}{https://github.com/LionSR/TNLean/issues/}
\newcommand{\tnleanPrBase}{https://github.com/LionSR/TNLean/pull/}
\newcommand{\tnleanDocBase}{https://sirui-lu.com/TNLean/docs/}
\newcommand{\ghissue}[1]{\href{\tnleanIssueBase#1}{GitHub issue \##1}}
\newcommand{\ghpr}[1]{\href{\tnleanPrBase#1}{PR \##1}}
\newcommand{\doclink}[1]{\href{\tnleanDocBase#1.html}{\path{#1}}}

% Dirac notation and the identity operator, as in blueprint/src/macros/common.tex.
% \providecommand keeps note-local definitions authoritative.
\usepackage{dsfont}
\providecommand{\ket}[1]{|#1\rangle}
\providecommand{\bra}[1]{\langle #1|}
\providecommand{\braket}[2]{\langle #1 | #2 \rangle}
\providecommand{\ketbra}[2]{|#1\rangle\!\langle #2|}
\providecommand{\Id}{\mathds{1}}
\providecommand{\one}{\mathds{1}}
\providecommand{\C}{\mathbb{C}}
\providecommand{\MN}[1]{M_{#1}(\C)}
\providecommand{\MD}{M_D(\C)}

% Verdict marker: \gapnote{<kind>}{<status>}. Kind is the policy
% classification (clarification, local-correction, scope-restriction,
% unfaithful, false-source, open-gap); status is open, wip (elimination
% underway), resolved, or historical. Machine-read by the site index;
% prints a verdict line.
\newcommand{\gapnote}[2]{\par\noindent\textbf{Verdict:} #1 (#2).\par}

\setlength{\emergencystretch}{2em}
\title{The Constant in the Bond-Dimension Bound\newline
for Repeated Qutrit States}
\author{}
\date{October 25, 2026}
\begin{document}
\maketitle
\gapnote{local-correction}{resolved}

\section{The printed bound}

Fix an integer $N\ge2$ and let $\psi$ be the unit qutrit state on $2N$ sites
with parameter $p_N=1/N$. Put $K=N^2$ and let $\chi=\psi^{\otimes K}$ be the
unit vector on $M=2N^3$ sites. For unit vectors $\eta,\phi$ write
$f=|\langle\eta,\phi\rangle|^2$ for the squared fidelity and
\begin{align}
    T(f) & =
    \begin{cases}
        2\sqrt{f(1-f)}, & f\ge1/2, \\
        1,              & f<1/2,
    \end{cases}
    \label{eq:swvc08_qutrit_optimized_error}
\end{align}
for the optimized pure-state trace error. Schuch, Wolf, Verstraete and
Cirac~\cite{Schuch2008EntropyScaling} define this error as the infimum of
$\lVert|\eta\rangle\langle\eta|-\alpha|\phi\rangle\langle\phi|\rVert_1$ over
scalars $\alpha$, which equals $\sin2\theta$ when $\cos\theta=\sqrt f$ and
$0\le\theta\le\pi/4$.\footnote{arXiv:0705.0292v2, \path{paper_v2.tex},
    lines 160--173.}
It is invariant under rescaling of $\phi$.
We write $T(\eta,\phi)$ for the value of
(\ref{eq:swvc08_qutrit_optimized_error}) at the squared fidelity of $\eta$
with $\phi/\lVert\phi\rVert$.

The argument of~\cite{Schuch2008EntropyScaling} combines two estimates. For
unit vectors $\eta,\phi$ and $K\ge1$, the first asserts
\begin{align}
    T(\eta^{\otimes K},\phi^{\otimes K})
     & \ge\sqrt{K/8}\,T(\eta,\phi)
    \qquad\text{whenever }T(\eta,\phi)^2\le2/K.
    \label{eq:swvc08_qutrit_printed_amplification}
\end{align}
The second asserts, for every matrix product state $\phi$ of bond dimension
at most $D$ on $2N$ sites,
\begin{align}
    T(\psi,\phi) & \ge\frac{2^N-(D-1)}{N2^N}.
    \label{eq:swvc08_qutrit_single_copy}
\end{align}
After reducing to a product approximant $\phi^{\otimes K}$, the source
concludes that every matrix product state $\Phi$ on $M$ sites of bond
dimension at most $D$ satisfies
\begin{align}
    D & \ge2^N\bigl(1-8T(\chi,\Phi)\bigr)+1.
    \label{eq:swvc08_qutrit_printed_bound}
\end{align}
The three displays occur at lines 471--479 of the source.

\section{What the printed estimates give}

Set $t=T(\chi,\Phi)$. When the hypothesis of
(\ref{eq:swvc08_qutrit_printed_amplification}) holds for the factor $\phi$,
combining it with (\ref{eq:swvc08_qutrit_single_copy}) and
$\sqrt{K/8}=N/(2\sqrt2)$ gives
\begin{align}
    \frac{2^N-(D-1)}{2\sqrt2\,2^N} & \le t,
    \label{eq:swvc08_qutrit_printed_combination} \\
    D                              & \ge2^N\bigl(1-2\sqrt2\,t\bigr)+1.
    \label{eq:swvc08_qutrit_sqrt_eight_bound}
\end{align}
Thus the printed constant eight is larger than the constant $2\sqrt2$ that
the printed ingredients yield. Since $t\ge0$,
(\ref{eq:swvc08_qutrit_sqrt_eight_bound}) implies
(\ref{eq:swvc08_qutrit_printed_bound}), so the printed bound is true but not
sharp for its argument. The argument also leaves the smallness hypothesis of
(\ref{eq:swvc08_qutrit_printed_amplification}) unchecked for the extracted
factor.

\section{The formal bound}

The formalization replaces (\ref{eq:swvc08_qutrit_printed_amplification}) by a
scalar estimate without a single-copy smallness hypothesis. For
$0\le f\le1$ and $K\ge1$,
\begin{align}
    T(f^K)<1 & \Longrightarrow\sqrt{K/4}\,T(f)\le T(f^K).
    \label{eq:swvc08_qutrit_formal_amplification}
\end{align}
Indeed, $T(f^K)<1$ forces $f\ge f^K>1/2$, and then
$Kf(1-f)\le4f^K(1-f^K)$, which is the square of
(\ref{eq:swvc08_qutrit_formal_amplification}).\footnote{%
    \leanid{Real.sqrt_mul_optimizedPureTraceError_le_pow_of_lt_one}, in the
    QICLean dependency.}
Squared fidelities of tensor products multiply, and the largest factor
fidelity $f_i$ of a product approximant satisfies $f_i^K\ge\prod_jf_j$.
Since $T$ is nonincreasing on $[1/2,1]$ and equal to one below $1/2$, this
gives $T(f_i^K)\le t$. A unit matrix product state $\Phi$ of bond dimension at
most $D$ on $M$ sites therefore yields a unit block state $\phi$ of bond
dimension at most $D$ on $2N$ sites with $(N/2)\,T(\psi,\phi)\le t$ whenever
$t<1$.\footnote{\leanid{TNLean.MPS.AreaLaw.IsOpenMPS.exists_copy_error_amplification}
    in \doclink{TNLean/MPS/AreaLaw/CopyErrorAmplification}.}
Combining this with (\ref{eq:swvc08_qutrit_single_copy}) gives
\begin{align}
    \frac{2^N-(D-1)}{2\cdot2^N} & \le t,
    \label{eq:swvc08_qutrit_formal_combination} \\
    D                           & \ge2^N(1-2t)+1.
    \label{eq:swvc08_qutrit_formal_bound}
\end{align}
For $t=1$, the right-hand side of (\ref{eq:swvc08_qutrit_formal_bound}) is
$1-2^N<0$, so (\ref{eq:swvc08_qutrit_formal_bound}) holds for every unit
approximant.\footnote{\leanid{TNLean.MPS.AreaLaw.EntropyScaling.bondDimension_qutritCopies_ge_of_scalar_error_lt_one}
    in \doclink{TNLean/MPS/AreaLaw/EntropyScaling/QutritApproximation}.}
It implies (\ref{eq:swvc08_qutrit_sqrt_eight_bound}) and
(\ref{eq:swvc08_qutrit_printed_bound}).

The approximation criterion of the source is the full trace-norm error
$\delta=\lVert|\chi\rangle\langle\chi|-|\Phi\rangle\langle\Phi|\rVert_1$ of
a possibly unnormalized vector $\Phi$.\footnote{arXiv:0705.0292v2,
    \path{paper_v2.tex}, lines 151--157.}
When $\Phi\ne0$, taking $\alpha=\lVert\Phi\rVert^2$ in the infimum that
defines $T(\chi,\Phi/\lVert\Phi\rVert)$ shows $T(\chi,\Phi)\le\delta$. If $\delta<1$, then $\Phi\ne0$, because the zero vector has
error one. Hence
\begin{align}
    D & \ge2^N(1-2\delta)+1
    \label{eq:swvc08_qutrit_full_error_bound}
\end{align}
for every matrix product state $\Phi$ of bond dimension at most $D$; for
$\delta\ge1$ the right-hand side is negative.\footnote{%
    \leanid{TNLean.MPS.AreaLaw.EntropyScaling.bondDimension_qutritCopies_ge}
    in \doclink{TNLean/MPS/AreaLaw/EntropyScaling/QutritApproximation}.}
The resulting obstruction to polynomial bond dimension is stated for
$\delta<1/8$, the threshold of (\ref{eq:swvc08_qutrit_printed_bound}), although
(\ref{eq:swvc08_qutrit_full_error_bound}) would permit any $\delta<1/2$.

\section{Verdict}

The printed bound (\ref{eq:swvc08_qutrit_printed_bound}) is true. Its
constant is not the one its printed ingredients give, and the smallness
hypothesis of (\ref{eq:swvc08_qutrit_printed_amplification}) is not verified
in the cited argument. The formalization proves the stronger bound
(\ref{eq:swvc08_qutrit_formal_bound}) by the amplification estimate
(\ref{eq:swvc08_qutrit_formal_amplification}), which needs only $t<1$. This is
a local correction of the constant and of the amplification step, with no
change to the hypotheses of the source.

\bibliographystyle{alpha}
\bibliography{references}
\end{document}
